\documentclass{beamer}
\usepackage{graphicx} % Required for inserting images
\usepackage{german}
\usepackage{amsmath}
\usepackage{amssymb}
\usepackage{amsthm}
\usepackage{url}


\newcommand{\R}{\mathbb{R}}% reele Zahlen
\newcommand{\C}{\mathbb{C}}% komplexe Zahlen
\newcommand{\Q}{\mathbb{Q}}% rationale Zahlen


\newtheorem{thm}{Satz}
\newtheorem{cor}[thm]{Korollar}
\newtheorem{lem}[thm]{Lemma}
\newtheorem{prop}[thm]{Proposition}

\title{Proseminar: Rekursion und Existenzaussagen}
\author{Valerie Höhle}
\date{\today}

\begin{document}

% Title Slide
\begin{frame}
    \titlepage
\end{frame}

\begin{frame}{Überblick}
\begin{enumerate}
    \item Rekursionen
    \item Diskrete Wahrscheinlichkeitsrechnung
    \item Existenzaussagen
\end{enumerate}
\end{frame}

\begin{frame}{Überblick}

\textbf{1. Rekursionen}
\begin{itemize}
    \item Binomialkoeffizienten
    \item Pascalsches Dreieck
    \item Stirling-Zahlen
\end{itemize}

\vspace{0.5cm}

2. Diskrete Wahrscheinlichkeitsrechnung

\vspace{0.2cm}

3. Existenzaussagen

\end{frame}


\begin{frame}{1. Rekursionen}

\textbf{Idee:}
Rekursionen ermöglichen es, neue Werte aus bereits bekannten kleineren Werten zu berechnen.

\vspace{0.4cm}

Zur Erinnerung betrachten wir die bereits bekannten
\textbf{Binomialkoeffizienten}:

[\binom{n}{k}
\frac{n!}{k!(n-k)!}]

[\binom{n}{k}
\binom{n}{n-k}]

[\binom{n}{k}
\frac{n(n-1)\cdots(n-k+1)}{k!}]

\end{frame}

\begin{frame}{1. Rekursionen}

\textbf{Erweiterung der Binomialkoeffizienten}

\begin{itemize}
    \item Wir wollen die Binomialkoeffizienten auf \(\mathbb{Z}\) und \(\mathbb{C}\) erweitern.
    
    \item Zunächst setzen wir
    \[
    \binom{0}{0}=1,
    \]
    da die leere Menge genau eine Teilmenge besitzt, nämlich sich selbst.

    \item Außerdem setzen wir
    \[
    0! = 1
    \]
    sowie
    \[
    n^{\underline{0}} = n^{\overline{0}} = 1.
    \]

    \item Für die Fakultät benötigen wir zunächst \(k \geq 0\).
\end{itemize}

\vspace{0.2cm}

\textbf{Definition 1.1 (Verallgemeinerte Binomialkoeffizienten)}

Für \(r\in\mathbb{C}\) und \(k\in\mathbb{Z}\) definieren wir

\[
\binom{r}{k}
=
\begin{cases}
\dfrac{r(r-1)\cdots(r-k+1)}{k!},
& k\geq 0,\\[1ex]
0,
& k<0.
\end{cases}
\]

\end{frame}

\end{document}
